\chapter{Stakeholder und Beduerfnisse}

\section{Stakeholderanalyse}

Stakeholder sind nicht nur direkte Nutzer. Fuer einen mobilen Roboter sind auch Betreiber, Wartung, Sicherheitsverantwortliche, IT, Entwicklung, Einkauf und Personen im Fahrbereich relevant. Das Lehrbuch hebt hervor, dass ein technisch funktionierendes System im Betrieb scheitern kann, wenn Wartbarkeit, Sicherheit, Datenschutz oder Betriebskontext nicht frueh beruecksichtigt werden \cite[Kap.~7]{becker2026sysml}.

\begin{longtable}{L{0.18\textwidth} L{0.30\textwidth} L{0.38\textwidth}}
\toprule
ID & Stakeholder & Interesse \\
\midrule
\modeltable{stakeholders_a1}
\bottomrule
\end{longtable}

\section{Stakeholder Needs}

\begin{longtable}{L{0.15\textwidth} L{0.18\textwidth} L{0.47\textwidth} L{0.10\textwidth}}
\toprule
ID & Quelle & Need & Prioritaet \\
\midrule
\modeltable{needs_a1_initial}
\bottomrule
\end{longtable}

\section{Welche Needs haben schon ein Requirement?}

Diese Tabelle ist keine Pflegeliste, sondern eine Antwort des Modells. Stakeholder sind dort Rollen (\texttt{part def}), Needs sind \texttt{concern}-Elemente mit ihren Stakeholdern, und jedes Requirement ist per Abhaengigkeit mit den Needs verbunden, aus denen es abgeleitet ist. Steht in der letzten Spalte \enquote{noch keins}, hat das Modell zu diesem Need noch keine Anforderung.

\begin{longtable}{L{0.18\textwidth} L{0.26\textwidth} L{0.42\textwidth}}
\toprule
Need & Stakeholder & Abgeleitete Requirements \\
\midrule
\modeltable{needs_coverage_a1}
\bottomrule
\end{longtable}

\section{Konflikte und Trade-offs}

\begin{longtable}{L{0.20\textwidth} L{0.32\textwidth} L{0.34\textwidth}}
\toprule
Konflikt & Beschreibung & Zu klaerende Entscheidung \\
\midrule
Sicherheit vs. Durchsatz & Niedrige Geschwindigkeit und konservatives Stoppen erhoehen Sicherheit, koennen aber Transportzeit verlaengern. & Zielgeschwindigkeit, Sicherheitsabstand, Zonenlogik \\
Kosten vs. Wahrnehmungsleistung & Hoeherwertige Sensorik verbessert Navigation und Hinderniserkennung, erhoeht aber Kosten und Energiebedarf. & Sensorvariante fuer MVP und Zielsystem \\
MVP-Fokus vs. Zielsystem & Ein kleiner MVP reduziert Aufwand, darf aber die spaetere Transportfunktion nicht architektonisch blockieren. & Welche Zielsystem-Schnittstellen werden bereits modelliert? \\
Kartierung vs. Stabilitaet & Automatische Kartenupdates koennen hilfreich sein, duerfen aber keine fehlerhaften dauerhaften Kartenveraenderungen erzeugen. & Trennung von globaler Karte, lokaler Karte und temporaeren Hindernissen \\
\bottomrule
\end{longtable}

\section{Detaillierte Stakeholder Needs aus Iteration A}

Die folgenden Needs wurden in der Vervollstaendigung von Iteration A feiner strukturiert. Sie ergaenzen die initiale Need-Liste und dienen als direkte Quelle fuer Iteration B.

\subsection{Bedienung und Auftrag}

\begin{longtable}{L{0.18\textwidth} L{0.16\textwidth} L{0.50\textwidth}}
\toprule
ID & Quelle & Need \\
\midrule
\modeltable{needs_a1_bedienung}
\bottomrule
\end{longtable}

\subsection{Transport}

\begin{longtable}{L{0.18\textwidth} L{0.16\textwidth} L{0.50\textwidth}}
\toprule
ID & Quelle & Need \\
\midrule
\modeltable{needs_a1_transport}
\bottomrule
\end{longtable}

\subsection{Navigation und Umgebung}

\begin{longtable}{L{0.18\textwidth} L{0.16\textwidth} L{0.50\textwidth}}
\toprule
ID & Quelle & Need \\
\midrule
\modeltable{needs_a1_navigation}
\bottomrule
\end{longtable}

\subsection{Safety}

\begin{longtable}{L{0.18\textwidth} L{0.16\textwidth} L{0.50\textwidth}}
\toprule
ID & Quelle & Need \\
\midrule
\modeltable{needs_a1_safety}
\bottomrule
\end{longtable}

\subsection{Betrieb, Wartung und Entwicklung}

\begin{longtable}{L{0.18\textwidth} L{0.16\textwidth} L{0.50\textwidth}}
\toprule
ID & Quelle & Need \\
\midrule
\modeltable{needs_a1_betrieb}
\bottomrule
\end{longtable}

\subsection{Unterricht und Demonstration}

\begin{longtable}{L{0.18\textwidth} L{0.16\textwidth} L{0.50\textwidth}}
\toprule
ID & Quelle & Need \\
\midrule
\modeltable{needs_a1_unterricht}
\bottomrule
\end{longtable}
